\documentclass[10pt, a4paper]{article}

% Parametri che modificano il file main.tex
% Le uniche parti da cambiare su main.tex sono:
% - vari \vspace tra sezioni
% - tabella azioni da intraprendere
% - sezione altro

\def\data{2024-03-12}
\def\oraInizio{14:00}
\def\oraFine{15:30}
\def\luogo{Piattaforma Discord}

\def\tipoVerb{Interno} % Interno - Esterno

\def\nomeResp{Orlandi G.} % Cognome N.
\def\nomeVer{Bresolin G.} % Cognome N.
\def\nomeSegr{Michelon R.} % Cognome N.

\def\nomeAzienda{Azzurro Digitale}
\def\firmaAzienda{azzurrodigitale.png}
\def\firmaResp{giacomo.png} % nome Responsabile

\def\listaPartInt{
Bresolin G.,
Campese M.,
Ciriolo I.,
Dugo A.,
Feltrin E.,
Michelon R.,
Orlandi G.
}

\def\listaPartEst{
Azzurro B.,
Digitale C.,
}

% Se nessuna revisione: \def\listaRevisioneAzioni {x}
\def\listaRevisioneAzioni {x}

\def\listaOrdineGiorno {
{Aggiornamento da parte di tutti i membri del gruppo riguardo al lavoro di codifica iniziato dopo la riunione precedente;},
{Preparazione per la riunione con l'azienda pianificata in data 2024-03-13;},
{Assegnazione dei task relativi alla parte di testing.}
}

\def\listaDiscussioneInterna {
{Il team ha eseguito un aggiornamento generale relativo alla codifica del backend in corso da una settimana. Considerati i buoni risultati il gruppo ha discusso della possibilità di presentare il lavoro svolto durante la riunione con AzzurroDigitale pianificata per il 2024-03-13;},
{Il team ha riportato la discussione tenuta in settimana per la divisione del lavoro da svolgere riguardante la parte di test di unità delle componenti sviluppate durante la prima parte di codifica del progetto. La suddivisione delle task viene quindi riportata in 'Azioni da intraprendere'.
}
}



% Se nessuna decisione: \def\listaDecisioni {x}
\def\listaDecisioni{
{Il team presenterà all'azienda il lavoro di codifica delle componenti di backend svolto fino ad ora;},
{Il team continuerà la produzione della parte di test di unità delle componenti di backend secondo la suddivisione dei compiti riportata in 'Azioni da intraprendere'.}
}

\usepackage{style}
\usepackage{headerfooter}

\title{\titolo}
\author{SWEetCode}

\begin{document}

% PRIMA PAGINA
\begin{titlepage}
    \thispagestyle{empty}
    \begin{tikzpicture}[remember picture, overlay]
        % TRIANGOLI
        \draw[fill=secondarycolor, secondarycolor] (current page.north west) -- (current page.south west) -- (8.8, -28);
        \draw[fill=primarycolor, primarycolor] (-3, 5) -- (4, -13.6) -- (11, 5);

        % LOGO
        \node [xshift=-5cm, yshift=25cm] (logo) at (current page.south east) {\includegraphics[width=6.5cm]{img/logo.png}};

        % SWEETCODE - DATE
        \node [anchor=north east, align=right, xshift=-1.2cm, yshift=20.5cm, text=black] (sweetcode) at (current page.south east) {\fontsize{32pt}{36pt}\selectfont SWEetCode};
        \draw[line width=4pt, lightcol] ([xshift=-3cm, yshift=-0.37cm]sweetcode.south west) -- ([yshift=-0.37cm]sweetcode.south east);
        \node [anchor=north east, align=right, xshift=-1.2cm, yshift=18.7cm, text=black] (date) at (current page.south east){\fontsize{24pt}{24pt} \selectfont Verbale \tipoVerb};

        % NOME FILE
        \node [anchor=north east, text width=15cm, align=right, xshift=-1.2cm, yshift=17cm, text=black] (titolo) at (current page.south east){\fontsize{48pt}{48pt}\textbf{\data}};

        % BOX DATI PARTECIPANTI
        \ifthenelse{\equal{\tipoVerb}{Esterno}}{
            \node[anchor=north east, xshift=-1.2cm, yshift=14.5cm, minimum width=8cm] (box) at (current page.south east){};
        }{
            \node[anchor=north east, xshift=-1.2cm, yshift=12.5cm, minimum width=8cm] (box) at (current page.south east){};
        }

        % RESPONSABILE
        \node[anchor=north west, align=left] (dati2) at (box.north west) {\fontsize{15pt}{15pt}\selectfont \textbf{Responsabile}};
        \draw[line width=4pt, lightcol] (dati2.south west) -- ([xshift=8cm]dati2.south west);
        \node[anchor=north west, align=left] (dati21) at (dati2.south west){\fontsize{13pt}{13pt}\selectfont \nomeResp};

        % VERIFICATORE
        \node[anchor=north west, yshift=-1cm, align=left] (dati3) at (dati21.north west) {\fontsize{15pt}{15pt}\selectfont \textbf{Verificatore}};
        \draw[line width=4pt, lightcol] (dati3.south west) -- ([xshift=8cm]dati3.south west);
        \node[anchor=north west, align=left] (dati31) at (dati3.south west){\fontsize{13pt}{13pt}\selectfont \nomeVer};

        % SEGRETARIO DI RIUNIONE
        \node[anchor=north west, yshift=-1cm, align=left] (dati4) at (dati31.north west) {\fontsize{15pt}{15pt}\selectfont \textbf{Segretario di Riunione}};
        \draw[line width=4pt, lightcol] (dati4.south west) -- ([xshift=8cm]dati4.south west);
        \node[anchor=north west, align=left] (dati41) at (dati4.south west){\fontsize{13pt}{13pt}\selectfont \nomeSegr};
        
        % UNIPD - SWE
        \node [xshift=4.4cm, yshift=2.3cm, draw, secondarycolor, text=white] (uni) at (current page.south west) {\fontsize{20pt}{20pt} \selectfont Università di Padova};
        \node [xshift=0.65cm, yshift=0.7cm, draw, secondarycolor, text=white, below=of uni] (corso) {\fontsize{20pt}{20pt}\selectfont Ingegneria del Software};

        % FIRMA AZIENDA
        %\ifthenelse{\equal{\tipoVerb}{Esterno}}{
           % \draw[line width=4pt, lightcol] ([xshift=-1.2cm, yshift=5.2cm]current page.south east) -- ([xshift=-8cm, yshift=5.2cm]current page.south east);
          %  \node [xshift=-4.8cm, yshift=6.3cm] (logo) at (current page.south east) {\includegraphics[width=6cm]{img/firme/\firmaAzienda}};
          %  \node[anchor=north west, xshift=12.9cm, yshift=4.95cm, align=left] at (current page.south west)
           % {\fontsize{13pt}{13pt}\selectfont L'azienda: \nomeAzienda};
       % }

       %  FIRMA
        \draw[line width=4pt, lightcol] ([xshift=-1.2cm, yshift=1.8cm]current page.south east) -- ([xshift=-8cm, yshift=1.8cm]current page.south east);
        \node [xshift=-4.8cm, yshift=2.45cm] (logo) at (current page.south east) {\includegraphics[width=6cm]{img/firme/\firmaResp}};
        \node[anchor=north west, xshift=12.9cm, yshift=1.45cm, align=left] at (current page.south west)
        {\fontsize{13pt}{13pt}\selectfont Il Responsabile: \nomeResp};
        
    \end{tikzpicture}
\end{titlepage}

% REGISTRO DELLE VERSIONI
%Registro in ordine dalla più recente alla meno recente!

{\renewcommand{\arraystretch}{1.5}
\section*{Registro delle versioni}

\begin{xltabular}{\textwidth}{c|c|c|c|X}
\label{tab:long}

\textbf{Versione} & \textbf{Data} & \quantities{\textbf{Responsabile di}\\\textbf{stesura}}& \textbf{Revisore} & \quantities{\textbf{Dettaglio e}\\\textbf{motivazioni}} \\
\endfirsthead

\textbf{Versione} & \textbf{Data} & \quantities{\textbf{Responsabile di}\\\textbf{stesura}}& \textbf{Revisore} & \quantities{\textbf{Dettaglio e}\\\textbf{motivazioni}} \\
\endhead

\multicolumn{5}{r}{{Continua nella pagina successiva}} \\
\endfoot

\endlastfoot

\hline
v2.23.0(1) & $2024-03-20$ & \quantities{Ciriolo I.} & Feltrin E. & Gestione e visualizzazione delle chat.\\
\hline
v2.21.0(1) & $2024-03-19$ & \quantities{Ciriolo I.} & Michelon R. & Prima stesura.\\
\hline
    
\end{xltabular}}
\newpage

% INDICE
\tableofcontents
\newpage

\listoffigures
\newpage
% INIZIO PAGINE
\section{Introduzione}
\subsection{Obiettivo del documento}
Lo scopo del documento è l'illustrazione delle funzionalità del prodotto \textit{Knowledge Management AI}. Il documento include: 
\begin{itemize}
    \item Come configurare prodotto;
    \item Descrizione dell'interfaccia utente del prodotto e le possibili interazioni con essa da parte degli utenti.
\end{itemize} 
Non viene descritta l'installazione del prodotto poichè nel contesto di questo progetto le figure di utente e di installatore/manutentore sono distinte, quindi non sarà l'utente a provvedere all'installazione dei software necessari all'uso dell'applicazione.
\subsection{Glossario}
Per evitare ambiguità e incomprensioni relative al linguaggio e ai termini utilizzati nella documentazione relativa al progetto viene presentato un Glossario. I termini ambigui o specifici presenti nello stesso, verranno identificati con un pedice |g|. All’interno dei documenti viene identificata con tale scrittura solo e soltanto la prima occorrenza presente nel testo di un termine definito nel Glossario.

\subsection{Riferimenti}
\subsubsection{Riferimenti normativi}
\begin{itemize}
    \item \textit{(Norme di progetto v2.0.0(0))};
    \item \textit{Regolamento del progetto didattico}: \\
        \url{https://www.math.unipd.it/~tullio/IS-1/2023/Dispense/PD2.pdf}\\
        (Ultimo accesso: 2024-03-14);
    \item \textit{Standard ISO/IEC 9126}:\\
        \url{https://it.wikipedia.org/wiki/ISO/IEC_9126}\\
        (Ultimo accesso: 2024-03-14).
    \end{itemize}
\subsubsection{Riferimenti informativi}
\begin{itemize}
        \item \textit{(Analisi dei requisiti v2.0.0(0))};
        \item \textit{Capitolato C1}: \textit{Knowledge Management AI.}
        \end{itemize}

\newpage
%\section{Installazione}
%\subsection{Prerequisiti}
%\subsection{Installazione del prodotto}
%\subsection{Accesso}

\section{Istruzioni per l'utente}
Questa sezione comprende le istruzioni per l'uso del prodotto delle funzionalità disponibili all'utente (attore).
\subsection{Configurazione utente}
La configurazione permette di selezionare diverse componenti del prodotto al primo avvio dell'applicazione come segue:
\subsubsection{Selezione del modello LLM}
\subsubsection{Selezione del vector database}
\subsubsection{Selezione del modello di embeddings}
\subsubsection{Selezione del sistema di archiviazione}

\subsection{Gestione dei documenti}
In questa sezione si illustrano le possibili azioni operabili sui documenti.
\subsubsection{Visualizzazione di un documento presente nel sistema di archiviazione}
Per visualizzare un documento presente nel sistema di archiviazione è necessario:
\begin{itemize}
    \item Selezionare dal menù in alto a sinistra nella dashboard la sezione 'File manager';
    \item Visualizzare la parte inferiore della pagina, sotto il titolo 'Lista dei documenti'. Alla destra di ogni documento della lista si visualizzano titolo, dimensioni, data di caricamento ed eventuale occultamento (stato);
    \item Selezionare con un clic il documento desiderato dalla lista.
\end{itemize}
\subsubsection{Rimozione dal sistema di archiviazione}
Per eliminare uno o più documenti presenti nel sistema di archiviazione è necessario:
\begin{itemize}
    \item Aprire il documento desiderato secondo i passaggi della sezione 2.2.1 'Visualizzazione di un documento presente nel sistema di archiviazione';
    \item Fare clic sul bottone 'Elimina documento' a destra del documento.
\end{itemize}

\subsubsection{Caricamento e rimozione dall'area di staging}
Per caricare uno o più documenti nell'area di staging è necessario:
\begin{itemize}
    \item Selezionare l'opzione 'File manager' nel menù in alto a sinistra della dashboard;
    \item Selezionare l'apposita icona per il caricamento in alto a destra della pagina, trascinandoci direttamente sopra il documento (o i documenti) da caricare oppure facendo un clic e selezionandoli dalla lista dei documenti personali;
    \item Se nell'area di staging è già presente un documento (o più documenti) con lo stesso nome e formato è possibile scegliere se sostituire il documento già presente nel sistema con quello nuovo cliccando su ..., oppure se non sostituirlo, selezionando l'opzione ... 
    \item  Verrà successivamente richiesta la conferma della scelta selezionata: per confermarla fare clic su ..., altimenti su ... per annullarla.
    \end{itemize} 
Per rimuovere uno o più documenti dall'area di staging è necessario ..


\subsubsection{Upload nel sistema di archiviazione e nel vector database}
Per poter fare l’upload nel sistema di archiviazione e nel vector database dei documenti caricati nell’area di staging è necessario ...

\subsubsection{Occultamento e riabilitazione dal sistema}
Per occultare uno o più documenti presenti nel sistema di archiviazione è necessario:
\begin{itemize}
    \item Aprire il documento desiderato secondo i passaggi della sezione 2.2.1 'Visualizzazione di un documento presente nel sistema di archiviazione';
    \item Fare clic sul bottone verde 'Occulta documento' alla destra del documento.
\end{itemize}
Per riabilitare uno o più documenti occultati è necessario:
\begin{itemize}
    \item Aprire il documento desiderato secondo i passaggi della sezione 2.2.1 'Visualizzazione di un documento presente nel sistema di archiviazione';
    \item Fare clic sul bottone verde 'Riabilita documento' alla destra del documento.
\end{itemize}

\subsubsection{Ricerca nel sistema}
Per cercare uno documento presente nel sistema di archiviazione è necessario:
\begin{itemize}
    \item Selezionare dal menù in alto a sinistra della dashboard l'opzione 'File manager';
    \item Fare clic sulla barra di ricerca (con scritto all'interno 'Cerca') situata in alto a destra della sezione 'Lista documenti', nella porzione inferiore della pagina;
    \item Digitare il titolo del documento ricercato.
\end{itemize}

\subsection{Gestione delle chat}
In questa sezione si illustrano le possibili azioni operabili sulle chat dell'applicazione. 
\subsubsection{Visualizzazione}
Per visualizzare la lista delle chat presenti nell’applicazione è sufficiente visualizzare la parte destra della dashboard dell'applicazione, sotto il titolo 'Lista ultime chat', dove è possibile visualizzare le singole chat presenti in lista automaticamente con il rispettivo titolo in ordine cronologico. Vengono visualizzati inoltre il timestamp dell'ultimo messaggio presente per ogni chat (nella rispettiva riga in alto a destra), la preview dell'ultimo messaggio presente della chat (appena sotto al titolo) e la preview dell'ultimo mittente di una chat.\\ Se si comincia una conversazione con il chatbot o se si seleziona l'opzione 'Chatbot' dal menù in alto a sinistra, la lista delle chat resta comunque disponibile nella parte sinistra della pagina e le sue funzionalità rimangono invariate.
\subsubsection{Creazione, eliminazione, rinominazione}
Per creare una nuova chat è necessario:
\begin{itemize}
    \item Selezionare il bottone 'Nuova chat' in alto a sinistra nella dashboard;
    \item In alternativa, selezionare dal menù in alto a sinistra l'opzione 'Chatbot' e selezionare il bottone 'Nuova chat' in alto a sinistra nella pagina che varrà aperta (pagina 'Chatbot').
\end{itemize}
Per eliminare una o più chat è necessario:
\begin{itemize}
    \item Selezionare dal menù in alto a sinistra della dashboard l'opzione 'Chatbot';
    \item Selezionare il simbolo dei tre puntini in verticale alla destra della chat che si vuole eliminare dalla lista delle chat (che compare nella parte sinistra della pagina). Se non si trova la chat desiderata, è possibile ricercarla seguendo la sezione 2.3.3 ('Ricerca');
    \item Selezionare l'opzione 'Elimina chat'.
\end{itemize}
Per rinominare una chat è necessario:
\begin{itemize}
    \item Selezionare dal menù in alto a sinistra della dashboard l'opzione 'Chatbot';
    \item Selezionare il simbolo dei tre puntini in verticale alla destra della chat che si vuole eliminare dalla lista delle chat (che compare nella parte sinistra della pagina). Se non si trova la chat desiderata, è possibile ricercarla seguendo la sezione 2.3.3 ('Ricerca');
    \item Selezionare l'opzione 'Rinomina chat';
    \item Inserire il nuovo titolo della chat.
\end{itemize}

\subsubsection{Ricerca}
Per cercare una chat o il contenuto di essa all'interno della lista delle chat è necessario:
\begin{itemize}
    \item Recarsi nella sezione 'Chatbot' nel menù in alto a sinistra della dashboard;
    \item Fare clic sull'apposita barra di ricerca con sopra la dicitura 'Cerca una chat' (e con all'interno la dicitura 'Cerca') situata nella parte sinistra della pagina, nella sezione 'Lista chat'.
\end{itemize} 
\subsubsection{Visualizzazione contenuto e messaggi}
Per visualizzare il contenuto ed i singoli messaggi di una chat è necessario:
\begin{itemize}
    \item Fare clic sulla chat che si intende leggere dalla lista delle chat. È possibile farlo sia dalla dashboard, selezionando la chat dalla lista delle ultime chat in alto a destra, o aprendo la pagina 'Chatbot' dal menù in alto a sinistra della dashboard e selezionandola dalla lista delle chat a sinistra nella pagina.
\end{itemize}
La chat desiderata verrà automaticamente aperta per la lettura. Ne vengono quindi visualizzati il contenuto testuale e l'orario dei singoli messaggi (in basso a sinistra rispetto ai messaggi). 
\subsubsection{Domande suggerite}
Se la chat è vuota vengono visualizzate (DOVE?) di default delle domande suggerite che possono essere poste al chatbot. Per selezionare una domanda è sufficiente selezionarla con un clic e confermarne l'invio selezionando l'opzione ...
\subsubsection{Inserimento, invio e copia}
Per inserire un messaggio nell'area di testo di una chat è necessario:
\begin{itemize}
    \item Recarsi nella sezione 'Chatbot' dal menù in alto a destra della dashboard;
    \item Digitare il testo del messaggio da inserire nell'apposito spazio in basso alla chat;
    \item Per inserirlo nella chat (quindi per inviarlo) si deve fare clic sul bottone azzurro a destra della casella per l'inserimento testuale.
\end{itemize}
Per poi copiare un messaggio è necessario fare clic sull'apposita icona con il simbolo di copia posta in alto a sinistra di ogni messaggio della chat.
\subsubsection{Modalità vocale}
È possibile utilizzare la modalità vocale, quindi inserire un messaggio in una chat tramite comando vocale, in questo modo:
\begin{itemize}
    \item Recarsi nella pagina 'Chatbot' dal menù in alto a sinistra della dashboard;
    \item Fare clic sul bottone del microfono a sinistra della casella di inserimento testuale;
    \item Scandire ad alta voce il testo del messaggio che si desidera inserire;
    \item Quando finito, premere nuovamente il bottone del microfono a sinistra della casella di inserimento testuale.
\end{itemize}
Per ascoltare un messaggio tramite audio è necessario ...





\end{document}